\documentclass[12pt,a4paper]{article}
\usepackage{amsmath,amssymb,amsthm}
\usepackage{graphicx}
\usepackage{geometry}
\geometry{a4paper, left=2.5cm, right=3cm, top=3cm, bottom=3cm}
\usepackage{fancyhdr}
\usepackage{setspace}
\usepackage{booktabs}
\usepackage{float}
\usepackage{hyperref}
\usepackage{xepersian}

\settextfont{X Nazanin}
\setdigitfont{X Nazanin}
\setlatintextfont{Times New Roman}

\makeatletter
\def\UseMathForPositioningText{}
\def\@kernel@tabular@init{}
\makeatother

\begin{document}

\begin{center}
{\Large\bfseries دانشگاه تهران}\\[0.2cm]
{\large\bfseries دانشکده مهندسی برق و کامپیوتر}\\[0.2cm]
{\large\bfseries گزارش فنی و تحلیل معماری مکانیزم توجه}\\[0.5cm]
\rule{\linewidth}{1pt}\\[0.3cm]
{\LARGE\bfseries ماژول توجه: cosformer}\\[0.2cm]
{\large \textbf{شناسه رسمی در فریم‌ورک:} \lr{Cosformer: Rethinking Softmax in Attention with Linear Complexity}}\\[0.2cm]
\rule{\linewidth}{1pt}\\[0.5cm]
\end{center}

\begin{figure}[H]
\centering
\framebox[0.9\textwidth]{
\begin{minipage}{0.85\textwidth}
\vspace{0.3cm}
\centering
\textbf{[محل الحاق دیاگرام تخصصی این مکانیزم]}\\[0.2cm]
\small \textbf{پرامپت مهندسی‌شده هوش مصنوعی جهت ساخت تصویر:}\\[0.15cm]
\footnotesize \lr{A detailed, high-tech academic architecture diagram of cosformer (Cosformer: Rethinking Softmax in Attention with Linear Complexity). Visualizing Query, Key, Value tensor transformations, matrix operations, and memory optimizations inside GPU hardware cache. Glowing cybernetic blueprint style, dark background, cyan and deep amber highlights.}
\vspace{0.3cm}
\end{minipage}
}
\caption{دیاگرام مفهومی و شماتیک جریان داده در مکانیزم توجه cosformer}
\end{figure}

\section*{۱. کارنامه تجربی و شاخص‌های بنچ‌مارک}
\begin{itemize}
\item \textbf{وضعیت ارزیابی در خط لوله:} \lr{PASS}
\item \textbf{دقت ثبت‌شده روی بنچ‌مارک:} 78.21 درصد
\item \textbf{امتیاز \lr{F1}:} 0.7894
\item \textbf{اوج مصرف حافظه فعال‌سازی (\lr{Peak VRAM}):} 337.54 مگابایت
\end{itemize}
\textbf{تحلیل و توجیه علمی نتایج:} این ماژول با ثبت دقت 78.21 درصد مصالحه مطلوبی میان سرعت اجرای استنتاج و بازنمایی ویژگی‌های توکن‌ها به نمایش می‌گذارد. تقریب‌های اعمال‌شده باعث مهار گلوگاه پهنای باند حافظه گردیده‌اند.

\section*{۲. مبانی تئوریک و انگیزه طراحی معماری}
Standard Multi-Head Self-Attention requires quadratic computational complexity mathcal O (L 2) and memory footprint mathcal O (L 2) due to the explicit computation of the Softmax attention matrix: [ operatorname Attn (Q, K, V) = operatorname Softmax left( frac Q K top sqrt d right) V ] **Cosformer** (Qin et al., 2022, ICLR) addresses this computational bottleneck by replacing Softmax with a **decomposed non-negative kernel mapping combined with a cosine-based positional reweighting mechanism**.

\section*{۳. روابط و فرمول‌بندی دقیق ریاضی}
2.1 Cosine Positional Kernel Decomposition

\section*{۴. تحلیل پیچیدگی زمانی و رفتار حافظه سخت‌افزاری}
---

\end{document}
